\section{Stationary-boundary twist support}
\label{sec:app_symmetry_stationary_twists}

Fix an MPS tensor $A=(A^i)_{i=0}^{d-1}$, a physical matrix $u$, and a
stationary boundary matrix $\Lambda\in\MN{D}$.  The twisted transfer map
$\E_u$ is Definition~\ref{def:twisted_transfer_map}.  The next constructions
are virtual-boundary functionals; they should not be confused with the
physical endpoint correlator of Definition~\ref{def:physical_string_order_param}.

\begin{definition}[Twisted companion tensor]\label{def:twisted_mixed_companion}
    \lean{MPSTensor.twistedMixedCompanion}
    \leanok
    \uses{def:mps_tensor}
    For an MPS tensor $A$ and a physical-index matrix $u$, define the
    companion tensor $B_u$ by
    \begin{align}
        B_u^n
        &:= \sum_{n'=0}^{d-1} \overline{u_{n'n}}\, A^{n'}.
        \label{eq:symmetry_companion_tensor}
    \end{align}
\end{definition}

\begin{lemma}[Twisted transfer as a mixed transfer operator]%
    \label{lem:twisted_transfer_eq_mixed_transfer}
    \lean{MPSTensor.twistedTransferMap_eq_mixedTransfer}
    \leanok
    \uses{def:twisted_transfer_map, def:twisted_mixed_companion, def:mixed_transfer}
    The twisted transfer map of $A$ is the mixed transfer operator of the pair
    $(A,B_u)$:
    \begin{align}
        \E_u &= F_{A,B_u}.
        \label{eq:symmetry_twisted_mixed_transfer}
    \end{align}
\end{lemma}

\begin{proof}
    \leanok
    Expand both definitions and rearrange the double sum.
\end{proof}


\begin{lemma}[The companion tensor has the same transfer map]
    \label{lem:transfer_map_twisted_companion}
    \lean{MPSTensor.transferMap_twistedMixedCompanion_eq}
    \leanok
    \uses{def:twisted_mixed_companion, thm:kraus_rectangular_freedom}
    If $uu^\dagger=\Id_d$, then the companion family $B_u$ satisfies
    \begin{align}
        \E_{B_u}(X)&=\E_A(X)
        \qquad (X\in\MN{D}).
        \notag
    \end{align}
\end{lemma}

\begin{proof}\leanok
    Expand $B_u^n=\sum_{n'}\overline{u_{n'n}}A^{n'}$.  The coefficient of
    $A^jX(A^k)^\dagger$ in $\E_{B_u}(X)$ is
    \begin{align}
        \sum_n\overline{u_{jn}}u_{kn}
        &=(uu^\dagger)_{kj}=\delta_{kj}.
        \notag
    \end{align}
    All off-diagonal terms vanish and the diagonal terms sum to $\E_A(X)$.
\end{proof}

\begin{definition}[Boundary string order parameter]%
    \label{def:string_order_boundary_param}
    \lean{MPSTensor.stringOrderBoundaryParam}
    \leanok
    \uses{def:twisted_transfer_map, def:string_order_param}
    For a boundary state $\Lambda$ and boundary matrices $X,Y \in \MN{D}$, the
    \emph{boundary string order parameter} is
    \begin{align}
        R_L(u;X,Y)
        &:= \tr\!(\Lambda X \E_u^L(Y)).
        \label{eq:symmetry_boundary_string}
    \end{align}
    The ordinary string order parameter is the special case $X=Y=\Id$.
    The matrices $X,Y$ act on the virtual level. In the string correlator
    of~\cite{PerezGarcia2008StringOrder} the string of twists is instead
    flanked by local operators $x,y$ acting on physical sites; expanding
    such operators in the transfer picture produces particular virtual
    matrices, and \cite{PerezGarcia2008StringOrder} argues that for
    injective tensors, physical operators on sufficiently many sites
    produce every virtual matrix. The results below are stated directly
    in this virtual-boundary form.
    % Deviation recorded in docs/paper-gaps/pgwsvc08_string_order_virtual_boundary.tex.
\end{definition}

\begin{lemma}[String order prevents decay to zero]
    \label{lem:not_tendsto_zero_of_hasStringOrder}
    \lean{MPSTensor.not_tendsto_zero_of_hasStringOrder}
    \leanok
    \uses{def:has_string_order}
    If $A$ has string order for $u$ with respect to a boundary state
    $\Lambda$, then for some boundary matrices $X,Y$ the sequence
    $R_L(u;X,Y)$ does not tend to $0$ as $L \to \infty$.
\end{lemma}

\begin{proof}
    \leanok
    \uses{def:has_string_order}
    A uniform lower bound $c \le \|R_L(u;X,Y)\|$ is incompatible with
    convergence to $0$.
\end{proof}


\section{Kraus mixing, transfer expansions, and a common TP gauge}
\label{sec:app_symmetry_common_tp_gauge}

The spectral estimates in the main chapter compare $A$ with $B_u$.  They are
placed in one trace-preserving gauge obtained from their common transfer map.

\begin{lemma}[Unitary Kraus mixing]\label{lem:unitary_kraus_mixing}
    \lean{MPSTensor.unitary_kraus_mixing}
    \leanok
    If $u$ is unitary, then replacing Kraus operators $\{A_i\}$ by
    their unitary mixtures $\{\sum_j u_{ij} A_j\}$ preserves the
    channel:
    \begin{align}
        \sum_i \left(\sum_j u_{ij} A_j\right)\, Y
        \left(\sum_j u_{ij} A_j\right)^\dagger
        &=
        \sum_i A_i\, Y\, A_i^\dagger.
        \label{eq:symmetry_unitary_kraus_mixing}
    \end{align}
\end{lemma}

\begin{proof}
    \leanok
    \uses{thm:kraus_same_map_of_unitary_combination}
    This is exactly the usual invariance of a Kraus decomposition under a
    unitary change of Kraus operators; see
    Theorem~\ref{thm:kraus_same_map_of_unitary_combination}.
\end{proof}


\begin{definition}[Common trace-preserving gauge for a twisted pair]
    \label{def:twisted_tp_gauge_setup}
    \lean{MPSTensor.TwistedTPGaugeSetup}
    \leanok
    \uses{def:twisted_mixed_companion}
    A \emph{common trace-preserving gauge setup} for $(A,u)$ consists of the
    companion $B_u$, a positive definite fixed point $\sigma$ of the adjoint
    transfer map, its positive square root $S=\sigma^{1/2}$, and the gauged
    tensors
    \begin{align}
        A'^i&=SA^iS^{-1},& B'^i&=SB_u^iS^{-1},
        \notag
    \end{align}
    together with
    $\sum_i(A'^i)^\dagger A'^i=\sum_i(B'^i)^\dagger B'^i=\Id$ and
    irreducibility of both gauged tensors, together with the inverse and
    adjoint identities for $S$ used in these calculations.
\end{definition}

\begin{theorem}[Construction of the common trace-preserving gauge]
    \label{thm:twisted_tp_gauge_setup}
    \lean{MPSTensor.twistedTPGaugeSetup}
    \leanok
    \uses{def:twisted_tp_gauge_setup, lem:transfer_map_twisted_companion,
        thm:psd_fp_unique}
    Let $D>0$.  If $A$ is injective, $uu^\dagger=\Id_d$, and
    $\E_A(\Id)=\Id$, then $(A,u)$ admits the setup of
    Definition~\ref{def:twisted_tp_gauge_setup}.
\end{theorem}

\begin{proof}\leanok
    Injectivity makes $\E_A$ irreducible.  The adjoint channel therefore has a
    positive definite stationary state $\sigma$.  By
    Lemma~\ref{lem:transfer_map_twisted_companion}, $A$ and $B_u$ have the same
    transfer map, hence the same adjoint fixed point.  Put $S=\sigma^{1/2}$.
    Positivity makes $S$ invertible, and direct substitution gives
    \begin{align}
       \sum_i(SA^iS^{-1})^\dagger(SA^iS^{-1})&=\Id,\\
       \sum_i(SB_u^iS^{-1})^\dagger(SB_u^iS^{-1})&=\Id.
       \notag
    \end{align}
    Similarity preserves irreducibility, so the two gauged families have all
    the asserted properties.
\end{proof}

\begin{lemma}[Twisted eigenvalues survive the common gauge]
    \label{lem:twisted_tp_gauge_has_eigenvalue}
    \lean{MPSTensor.twistedTPGaugeSetup_hasEigenvalue}
    \leanok
    \uses{def:twisted_tp_gauge_setup,
        lem:twisted_transfer_eq_mixed_transfer}
    If $V\ne0$ and $\E_u(V)=\lambda V$, then $\lambda$ is an eigenvalue of the
    mixed transfer map $F_{A',B'}$ in any common trace-preserving gauge setup.
\end{lemma}

\begin{proof}\leanok
    By Lemma~\ref{lem:twisted_transfer_eq_mixed_transfer}, the original
    equation is $F_{A,B_u}(V)=\lambda V$.  Multiplying on the left and right by
    $S$ and using $S^\dagger=S$ gives
    \begin{align}
       F_{A',B'}(SVS)&=\lambda SVS.
       \notag
    \end{align}
    Since $S$ is invertible, $SVS\ne0$.
\end{proof}

\section{Boundary and limit support}
\label{sec:app_symmetry_boundary_limits}

\begin{lemma}[Spectral radius $<1$ forces boundary-string decay]%
    \label{lem:boundary_string_order_tendsto_zero}
    \lean{MPSTensor.stringOrderBoundaryParam_tendsto_zero_of_spectralRadius_lt_one}
    \leanok
    \uses{def:string_order_boundary_param, def:twisted_transfer_map}
    If $\rho(\E_u) < 1$, then, as $L \to \infty$, for all boundary matrices
    $X,Y,\Lambda \in \MN{D}$ one has
    $R_L(u;X,Y) \longrightarrow 0$.
\end{lemma}

\begin{proof}
    \leanok
    In finite dimension, $\rho(\E_u) < 1$ implies $\E_u^L \to 0$ in operator
    norm. Composing with the fixed trace pairing
    $Z \mapsto \tr(\Lambda X Z)$ gives the stated scalar convergence.
\end{proof}

\section{Routine universality corollary}
\label{sec:app_symmetry_universality}

The following consequence uses only the virtual MPV-family symmetry predicate
and the canonical fixed-point hypotheses from the main chapter.

\begin{theorem}[Virtual-boundary nondecay agrees for two injective MPV-symmetric tensors]%
    \label{thm:has_string_order_iff_of_symmetric_injective}
    \lean{MPSTensor.hasStringOrder_iff_of_symmetric_injective}
    \leanok
    \uses{def:has_string_order, def:on_site_symmetric, def:injective}
    If $A$ and $B$ are injective, both on-site symmetric under a
    unitary representation $U$, and if $\E_A(\Id)=\Id$,
    $\E_B(\Id)=\Id$, $\Lambda_A$ and $\Lambda_B$ are positive
    definite, $\tr(\Lambda_A)=1$, $\tr(\Lambda_B)=1$,
    $\E_A^\dagger(\Lambda_A)=\Lambda_A$, and
    $\E_B^\dagger(\Lambda_B)=\Lambda_B$, then for every $g \in G$,
    string order exists for $A$ with twist $U(g)$ if and only if it
    exists for~$B$. In particular, since
    Definition~\ref{def:is_same_spt_phase} implies on-site symmetry
    for both tensors, this theorem applies to tensors in the same SPT
    phase; it applies equally to symmetric tensors in different SPT
    phases.
\end{theorem}

\begin{proof}\leanok
    \uses{thm:has_string_order_of_symmetric_injective}
    Both directions follow from
    Theorem~\ref{thm:has_string_order_of_symmetric_injective}, which
    shows that string order holds universally for injective symmetric
    tensors with canonical normalisations.
\end{proof}
